% Part: model-theory
% Chapter: models-of-arithmetic
% Section: non-standard-models

\documentclass[../../../include/open-logic-section]{subfiles}

\begin{document}

\olfileid{mod}{mar}{mdq}
\section{$\Th{Q}$ ची प्रतिमाने}

\begin{explain}
$\Struct{N}$ जी !!{sentence}s सत्य करते तीच सत्य करणाऱ्या अमानक
!!{structure}s आहेत, म्हणजे त्या~$\Th{TA}$ ची प्रतिमाने आहेत, हे आपल्याला
माहीत आहे. $\Sat{N}{\Th{Q}}$ असल्यामुळे $\Th{TA}$ चे कोणतेही प्रतिमान
$\Th{Q}$ चेही प्रतिमान असते. $\Th{Q}$ ही~$\Th{TA}$ पेक्षा फारच दुर्बल
आहे; उदाहरणार्थ, $\Th{Q} \Proves/ \lforall[x][\lforall[y][\eq[(x +
      y)][(y+x)]]]$. दुर्बलतर उपपत्ती पूर्ण करणे सोपे असते: त्यांना अधिक
प्रतिमाने असतात. उदाहरणार्थ, $\Th{Q}$ च्या काही प्रतिमानांत
$\lforall[x][\lforall[y][\eq[(x + y)][(y+x)]]]$ असत्य असते; पण ती
$\Th{TA}$ ची किंवा त्या दृष्टीने $\Th{PA}$ चीसुद्धा प्रतिमाने असू शकत
नाहीत. $\Th{Q}$ ची प्रतिमाने तुलनेने सोपीही असतात: ती स्पष्टपणे निश्चित
करता येतात.
\end{explain}

\begin{ex}
\ollabel{ex:model-K-of-Q}
$\Domain{K} = \Nat \cup \{a\}$ असा प्रांत आणि पुढील अर्थनिर्धारणे असलेली
!!{structure}~$\Struct{K}$ विचारात घ्या:
\begin{align*}
  \Assign{\Obj{0}}{K} & = 0\\
  \Assign{\prime}{K}(x) & =
  \begin{cases}
    x+1 & \text{जर $x\in \Nat$}\\
    a & \text{जर $x = a$}
  \end{cases}\\
  \Assign{+}{K}(x, y) & =
  \begin{cases}
    x+y & \text{जर $x$, $y \in\Nat$}\\
    a & \text{इतर प्रसंगी}
  \end{cases}\\
  \Assign{\times}{K}(x, y) & =
  \begin{cases}
    xy & \text{जर $x$, $y \in\Nat$}\\
    0 & \text{जर $x = 0$ किंवा $y = 0$}\\
    a & \text{इतर प्रसंगी}\\
  \end{cases}\\
  \Assign{<}{K} & =
  \Setabs{\tuple{x,y}}{x, y \in \Nat \text{ आणि } x<y} \cup
  \Setabs{\tuple{x,a}}{x \in \Domain{K}}
\end{align*}
$\Sat{K}{\Th{Q}}$ हे दाखवण्यासाठी $\Th{Q}$ ची सर्व स्वयंसिद्धके
$\Struct{K}$ मध्ये सत्य आहेत हे पडताळावे लागेल. सोयीसाठी,
$\Assign{\prime}{K}(x)$ करिता $x^\nssucc$ ($\Struct{K}$ मधील $x$ चा
``उत्तरवर्ती''), $\Assign{+}{K}(x, y)$ करिता $x \nsplus y$
($\Struct{K}$ मधील $x$ आणि $y$ यांची ``बेरीज''),
$\Assign{\times}{K}(x, y)$ करिता $x \nstimes y$ ($\Struct{K}$ मधील
$x$ आणि~$y$ यांचा ``गुणाकार''), आणि $\tuple{x,y} \in \Assign{<}{K}$
करिता $x \nsless y$ लिहू. या संक्षेपांनी $\Struct{K}$ मधील क्रिया अधिक
स्पष्टपणे अशा देता येतात:
\[
\begin{array}{c|c}
  x & x^\nssucc \\
  \hline
  n & n+1 \\
  a & a
\end{array}
\qquad
\begin{array}{c|ccc}
  x \nsplus y & 0 & m & a \\
  \hline
  0 & 0 & m & a \\
  n & n & n+m & a \\
  a & a & a & a \\
\end{array}
\qquad
\begin{array}{c|ccc}
  x \nstimes y & 0 & m & a \\
  \hline
  0 & 0 & 0 & 0 \\
  n & 0 & nm & a \\
  a & 0 & a & a \\
\end{array}
\]
$n$, $m \in \Nat$ यांच्यासाठी $n \nsless m$ हे तेव्हाच, जेव्हा $n<m$;
आणि प्रत्येक~$x \in \Domain{K}$ साठी $x \nsless a$.

$\nssucc$ हे !!{injective} असल्यामुळे
$\Sat{K}{\lforall[x][\lforall[y][(\eq[x'][y'] \lif \eq[x][y])]]}$.
$\Struct{K}$ मध्ये $0$ हा $\nssucc$-उत्तरवर्ती नसल्यामुळे
$\Sat{K}{\lforall[x][\eq/[\Obj 0][x']]}$. प्रत्येक $n>0$ साठी
$n = (n-1)^\nssucc$ आणि $a = a^\nssucc$ असल्यामुळे
$\Sat{K}{\lforall[x][(\eq[x][\Obj 0] \lor \lexists[y][\eq[x][y']])]}$.

$n \nsplus 0 = n+0 = n$ आणि~$\nsplus$ च्या व्याख्येवरून
$a\nsplus 0 = a$ असल्यामुळे $\Sat{K}{\lforall[x][\eq[(x + \Obj 0)][x]]}$.
$\Sat{K}{\lforall[x][\lforall[y][\eq[(x + y')][(x + y)']]]}$ हे थोडे
अधिक गुंतागुंतीचे आहे. $n$, $m$ या दोन्ही मानक असतील, तर:
\begin{align*}
(n \nsplus m^\nssucc) & = (n+(m+1)) = (n+m)+1 = (n \nsplus m)^\nssucc
\intertext{कारण मानक संख्यांवर $\nsplus$ आणि $^\nssucc$ ही अनुक्रमे $+$
  आणि $\prime$ यांच्याशी जुळतात. आता $x \in \Domain{K}$ असे समजा. मग}
(x \nsplus a^\nssucc) & = (x \nsplus a) = a = a^\nssucc = (x \nsplus a)^\nssucc
\intertext{उरलेले प्रकरण $y \in \Domain{K}$ पण $x =
  a$ असण्याचे आहे. येथे $y = n$ मानक आहे की $y = a$, यानुसारही प्रकरणे
  वेगळी करावी लागतात:}
(a \nsplus n^\nssucc) & = (a \nsplus (n+1)) = a = a^\nssucc = (a \nsplus n)^\nssucc\\
(a \nsplus a^\nssucc) & = (a \nsplus a) = a = a^\nssucc = (a \nsplus a)^\nssucc
\end{align*}
%READERNOTE{OLFOL-034 — गोठवलेल्या इंग्रजीतील प्रकरण-विभागणीत K च्या प्रांताबाहेरील b लिहिले आहे; पण लगेचची गणना y = a हेच प्रकरण हाताळते. मराठी स्पष्टीकरणात b ऐवजी a केले आहे.}
अर्थात, हा तपशील आवश्यकतेपेक्षा थोडा अधिक आहे. उदाहरणार्थ,
$z$ काहीही असले तरी $a \nsplus z = a$ असल्यामुळे
$a \nsplus a^\nssucc = a$ हे लगेच निष्पन्न होते. उरलेली स्वयंसिद्धकेही
याच रीतीने पडताळता येतात.

त्यामुळे $\Struct{K}$ हे~$\Th{Q}$ चे प्रतिमान आहे. तिची ``बेरीज''~$\nsplus$
ही क्रमनिरपेक्षसुद्धा आहे. पण $\Struct{N}$ मध्ये सत्य आणि $\Struct{K}$
मध्ये असत्य अशी इतर वाक्ये आहेत, तसेच याच्या उलटही आहे. उदाहरणार्थ,
$a \nsless a$, म्हणून $\Sat{K}{\lexists[x][x < x]}$ आणि
$\Sat/{K}{\lforall[x][\lnot x<x]}$. यावरून $\Th{Q} \Proves/
\lforall[x][\lnot x < x]$ हे दिसते.
\end{ex}

\begin{prob}
\olref[mod][mar][mdq]{ex:model-K-of-Q} मधील $\Struct{K}$ ही~$\Th{Q}$ ची
उरलेली स्वयंसिद्धके पूर्ण करते हे सिद्ध करा:
\begin{align*}
  & \lforall[x][\eq[(x \times \Obj 0)][\Obj 0]] \tag{$!Q_6$}\\
  & \lforall[x][\lforall[y][\eq[(x \times y')][((x \times y) + x)]]] \tag{$!Q_7$}\\
  & \lforall[x][\lforall[y][(x < y \liff \lexists[z][\eq[(z' + x)][y]])]] \tag{$!Q_8$}
\end{align*}
फक्त~$\prime$ असलेले, $\Struct{N}$ मध्ये सत्य पण $\Struct{K}$ मध्ये असत्य
असे !!a{sentence} शोधा.
\end{prob}

\begin{ex}
\ollabel{ex:model-L-of-Q} $\Domain{L} = \Nat \cup \{a, b\}$ असा प्रांत
आणि पुढीलप्रमाणे दिलेली $\Assign{\prime}{L} = \nssucc$ व
$\Assign{+}{L} = \nsplus$ ही अर्थनिर्धारणे असलेली !!{structure}~$\Struct{L}$
विचारात घ्या:
\[
\begin{array}{c|c}
  x & x^\nssucc \\
  \hline
  n & n+1 \\
  a & a\\
  b & b
\end{array}
\qquad
\begin{array}{c|ccc}
  x \nsplus y & m & a & b\\
  \hline
  n & n+m & b & a\\
  a & a & b & a\\
  b & b & b & a
\end{array}
\]
$\nssucc$ हे !!{injective} आहे, $0$ त्याच्या व्याप्तीत नाही, आणि
$0$ व्यतिरिक्त प्रत्येक~$x \in \Domain{L}$ त्याच्या व्याप्तीत आहे;
म्हणून $!Q_1$--$!Q_3$ ही स्वयंसिद्धके~$\Struct{L}$ मध्ये सत्य आहेत.
प्रत्येक $x$ साठी $x \nsplus 0 = x$, म्हणून $!Q_4$ सुद्धा सत्य आहे.
$!Q_5$ साठी $x \nsplus y^\nssucc$ आणि $(x \nsplus y)^\nssucc$ विचारात
घ्या. $x$ आणि $y$ दोन्ही मानक असतील, तर ते समान आहेत; कारण तेव्हा
$\nssucc$ आणि $\nsplus$ ही $\prime$ आणि $+$ यांच्याशी जुळतात. $x$ अमानक
आणि $y$ मानक असेल, तर $x \nsplus y^\nssucc = x = x^\nssucc =
(x \nsplus y)^\nssucc$. $x$ आणि $y$ दोन्ही अमानक असतील, तर चार प्रकरणे
आहेत:
\begin{align*}
&  a \nsplus a^\nssucc  = b = b^\nssucc = (a \nsplus a)^\nssucc\\
&  b \nsplus b^\nssucc  = a = a^\nssucc = (b \nsplus b)^\nssucc\\
&  b \nsplus a^\nssucc  = b = b^\nssucc = (b \nsplus a)^\nssucc\\
&  a \nsplus b^\nssucc  = a = a^\nssucc = (a \nsplus b)^\nssucc\\
\intertext{जर $x$ मानक पण $y$ अमानक असेल, तर}
&  n \nsplus a^\nssucc  = n \nsplus a = b = b^\nssucc = (n \nsplus a)^\nssucc\\
&  n \nsplus b^\nssucc  = n \nsplus b = a = a^\nssucc = (n \nsplus b)^\nssucc
\end{align*}
%READERNOTE{OLFOL-035 — गोठवलेल्या इंग्रजीतील तिसऱ्या गणनेच्या शेवटी मुक्त y आहे; त्या ओळीची डावी बाजू आणि L ची बेरीज-सारणी दोन्ही येथे a आवश्यक ठरवतात. मराठी सूत्रात y ऐवजी a केले आहे.}
त्यामुळे $\Sat{L}{!Q_5}$. पण $a \nsplus 0 \neq 0 \nsplus a$, म्हणून
$\Sat/{L}{\lforall[x][\lforall[y][\eq[(x+y)][(y+x)]]]}$.
\end{ex}

\begin{prob}
\olref[mod][mar][mdq]{ex:model-L-of-Q} मधील $\Struct{L}$ मध्ये
$\times$ आणि $<$ यांचे अर्थनिर्धारण करणारी $\nstimes$ आणि $\nsless$
समाविष्ट करून तिचा विस्तार करा. तुमची रचना~$\Th{Q}$ ची उरलेली
स्वयंसिद्धके पूर्ण करते हे दाखवा:
\begin{align*}
& \lforall[x][\eq[(x \times \Obj 0)][\Obj 0]] \tag{$!Q_6$}\\ &
  \lforall[x][\lforall[y][\eq[(x \times y')][((x \times y) + x)]]]
  \tag{$!Q_7$}\\ & \lforall[x][\lforall[y][(x < y \liff \lexists[z][\eq[(z'+x)][y]])]] \tag{$!Q_8$}
\end{align*}
\end{prob}

\begin{prob}
\olref[mod][mar][mdq]{ex:model-L-of-Q} मधील $\Struct{L}$ मध्ये $a^\nssucc
= a$ आणि $b^\nssucc = b$. ज्या~$\Th{Q}$ च्या प्रतिमानात
$a^\nssucc = b$ आणि $b^\nssucc = a$ असे आहे, असे प्रतिमान अस्तित्वात आहे का?
\end{prob}

\begin{explain}
अमानक !!{element}s हे मानक घटकांच्या ``पलीकडे'' असलेली~$\Th{Q}$ ची
प्रतिमाने आपण स्पष्टपणे रचली आहेत. किंबहुना, स्वयंसिद्धकांनुसार तितके असणे
आवश्यक आहे. अमानक !!{element}~$x$ हा ${} \nsless 0$ असू शकत नाही, कारण
$\Th{Q} \Proves \lforall[x][\lnot x<0]$
(\olref[inc][req][min]{lem:less-zero} पाहा). तसेच, प्रत्येक $n$ साठी
$\Th{Q} \Proves \lforall[x][(x < \num{n}' \lif
(\eq[x][\num{0}] \lor \eq[x][\num{1}] \lor \dots \lor
\eq[x][\num{n}]))]$ (\olref[inc][req][min]{lem:less-nsucc}); म्हणून
कोणत्याही~$n>0$ साठी $a \nsless n$ असू शकत नाही.
\end{explain}

\end{document}
